\section{State of the Art}
\label{state-of-the-art}
This section reviews prior work related to the main contributions of this study, covering the availability of Zigbee network traffic datasets, existing approaches to Zigbee traffic analysis and profiling, and prior efforts on the trade-off between traffic storage requirements and analysis accuracy, which have so far been explored primarily in Wi-Fi and IP-based network settings.

\subsection{Dataset Availability}
Recent survey works provide a comprehensive overview of publicly available network traffic datasets for IoT systems, consistently highlighting the severe scarcity of datasets containing realistic Zigbee traffic. In particular, out of 74 datasets surveyed in \cite{de2023survey}, only seven include Zigbee traffic, and only a small subset of these provide packet-level captures suitable for fine-grained traffic analysis, such as device fingerprinting or physical-layer studies.

Among the publicly available datasets, the work in \cite{c7-nvc6-4q28-22} is the only one that includes traffic from ten commercial Zigbee devices commonly deployed in smart home environments. Traffic was captured using a software-defined radio platform (USRP N210) with GNU Radio and the gr-ieee802-15-4 and gr-foo modules, enabling IEEE 802.15.4 packet reception and storage in PCAP format. However, the dataset is released exclusively as unlabeled raw traffic traces, limiting its applicability to supervised IoT forensic tasks.

Most other studies rely on proprietary or non-public datasets. For instance, Sadikin et al. \cite{sadikin2020zigbee} developed a Zigbee intrusion detection system and collected attack traces involving one gateway and three smart light bulbs, but did not release the dataset. Bihl et al. \cite{bihl2016feature} conducted physical-layer fingerprinting experiments on ten identical Texas Instruments CC2420 Zigbee devices by collecting RF emissions in controlled environments. Patel et al. \cite{patel2014improving} extended this line of work to 13 CC2420 devices and evaluated ensemble learning techniques, such as Random Forest and AdaBoost, achieving improved identification performance compared to parametric classifiers. Despite their promising results, the datasets used in these studies are not publicly available.

Among the few public datasets that include Zigbee traffic, Dadkhah et al. \cite{iot_dataset_paper} introduced a multidimensional IoT profiling dataset primarily focused on IP-based devices, with a limited Zigbee subset. This subset includes nine distinct Zigbee devices (16 instances) from four manufacturers, captured under controlled laboratory conditions. While noteworthy for incorporating non-IP IoT technologies, the Zigbee component of the dataset exhibits limited interaction diversity, as only single, non-repetitive device–action pairs were recorded. Such constraints limit its suitability for comprehensive traffic profiling in complex and distributed environments such as smart homes.
Lourme et al. \cite{NDW74U_2023} and Keleşoğlu et al. \cite{app142310844} released ZBDS2023 and ZigBeeNet, respectively: multi-day Zigbee traffic traces collected from 10 and 14 devices, with RSSI information available. ZBDS2023 further provides multi-point captures using multiple sniffers deployed at different locations within the home. Despite these strengths, both datasets offer limited support for evaluating generalization in realistic smart-home settings, as the devices are largely drawn from a single vendor ecosystem (predominantly Philips) and the captures do not span multiple network topologies, preventing robustness assessments of ML models across diverse deployment conditions.


Motivated by these limitations, the present work introduces a dedicated Zigbee data acquisition campaign aimed at increasing both interaction diversity and traffic volume. In contrast to prior datasets, our campaign captures traffic from 21 Zigbee devices spanning multiple functional classes, including temperature sensors, door sensors, motion detectors, smart plugs, smart bulbs, vibration sensors, and light sensors, and collects data in multi-topology environments where devices operate concurrently rather than in isolation. This setup enables a systematic analysis of how network topology and deployment conditions affect the performance and generalization of state-of-the-art traffic-profiling methodologies.

% approach captures device interactions within multi-topology network configurations, enabling the investigation of how network topology impacts state-of-the-art traffic analysis methodologies.

\subsection{Zigbee Network Traffic Analysis}
Despite the extensive literature on IoT traffic classification, relatively few studies focus specifically on Zigbee-based traffic analysis. Rong Li et al. \cite{ZPA2023} investigated potential privacy leakage from encrypted Zigbee traffic, aiming to infer device type and operational state. Their approach extends classical feature-based traffic analysis by incorporating inter-event timing information, along with selected protocol-level features such as MAC addresses and channel identifiers. Machine learning models were then employed for device classification, identification, and event detection.

Similarly, the authors of \cite{ziot2020} proposed a fingerprinting framework for the passive identification of device classes in Zigbee and Z-Wave networks. Their method primarily relies on inter-arrival time (IAT) statistics of network frames, achieving effective device-class identification. Other works, such as \cite{Peek-a-Boo2018,antonio}, also explored Zigbee traffic analysis, but were limited by the use of small-scale datasets.

Additional studies have examined encrypted Zigbee traffic from different perspectives. Some focus on exploiting protocol vulnerabilities to infer sensitive device information, such as logical device roles \cite{zigator20} or spatial location \cite{Localize2022}. Others have released Zigbee datasets targeting either device profiling \cite{iot_dataset_paper} or security applications, such as the dataset introduced in \cite{ZBDS2023} for intrusion detection. In parallel, several works have addressed Zigbee device fingerprinting using physical-layer characteristics rather than network traffic features \cite{27patel2015non,28bihl2016feature,30dubendorfer2013zigbee}.

While existing studies demonstrate the feasibility of Zigbee traffic classification, they generally overlook the impact of mesh network topology on model robustness. Unlike IP-based networks, Zigbee extensively relies on multi-hop communication, where changes in network topology may significantly affect observed traffic patterns. In this work, we build upon methodologies proposed in the state of the art to systematically quantify performance degradation across different network topologies, providing a baseline for machine-learning-based Zigbee traffic profiling in realistic large-scale smart home environments.

\subsection{Storage-accuracy analysis}
The trade-off between traffic storage requirements and the accuracy of downstream forensic analysis has received limited attention in IoT forensics, where most approaches assume the availability of complete traffic traces and overlook long-term storage scalability issues.
An initial study of this problem was presented in \cite{fabio2,fabio3}, which proposed a framework for optimizing storage efficiency while preserving analysis accuracy. Focusing on Wi-Fi and IP-based IoT traffic, the authors compared lossless compression of raw traces with lossy compression of extracted traffic features, showing that feature-level compression can substantially reduce storage requirements without degrading classification performance.
More recently, Palmese et al. \cite{palmese2025resource} investigated resource optimization for evidence collection and preservation in forensics-ready access points, analyzing how storage and processing constraints affect IP-based traffic monitoring at the network edge.
Despite these contributions, existing works remain limited to Wi-Fi and IP-based technologies and do not consider the distinctive characteristics of Zigbee networks, such as low traffic rates, small packet sizes, and mesh-based multi-hop communication. This work extends the storage--accuracy analysis to Zigbee traffic, evaluating compression strategies in realistic multi-topology smart home scenarios.